\section{Actual affine prime selectors and complete signed compensation}
\label{sec:affine-owner-selectors}

This section gives an additive construction in which different primes
may select different actual input roots. Its output is a new primitive
Eisenstein integer. In particular, preserving the selected prime powers
does not preserve the input power representation. All data in the
positive-density theorem below are fixed before the parameter tends
to infinity.

Write
\[
 \mathcal O=\mathbb Z[\zeta],\qquad \zeta^2-\zeta+1=0,
 \qquad N(a+b\zeta)=a^2+ab+b^2,
\]
and identify an element with its pair of integer coordinates. Put
\[
 P(a+b\zeta)=ab(a+b),\qquad |Z|=\sqrt{N(Z)}.
\]
For nonzero \(Z\), its content is the positive gcd of its coordinates.
When \(P(Z)\ne0\), the complete signed boundary cost is
\begin{equation}
 J(Z)=\sum_{p\mid |P(Z)|}\bigl(v_p(P(Z))-3\bigr)\log p
     =\log|P(Z)|-3\log\operatorname{rad}(|P(Z)|).
 \label{lm-signed-cost}
\end{equation}
In particular, depth one and depth two contribute respectively
\(-2\log p\) and \(-\log p\). These terms are retained throughout.

\begin{proposition}[The actual additive quotient and a short kernel]
\label{lm-additive-kernel}
For a finite set \(I\) and a positive integer \(n\), let
\[
 S:\mathbb Z^I\longrightarrow\mathcal O,\qquad
 S(c)=\sum_{k\in I}c_k w_k^n,\qquad w_k=3k+\zeta.
\]
Its integer image has rank at most two. Every family of homogeneous
congruences depending only on \(S(c)\), such as
\(\ell_q(S(c))\equiv0\pmod {q^{e_q}}\) for integer linear forms
\(\ell_q\), contains \(\ker S\). On any \(n+2\) consecutive indices,
there is the actual identity
\begin{equation}
 \sum_{j=0}^{n+1}(-1)^j\binom{n+1}{j}w_{k+j}^{n}=0.
 \label{lm-finite-difference}
\end{equation}
Its coefficient vector is nonzero and every coefficient has absolute
value at most \(2^{n+1}\).
\end{proposition}
\begin{proof}
Both coordinates of \((3X+\zeta)^n\) are integer polynomials of degree
at most \(n\). For a monomial of degree \(m\le n\), apply
\((T\,d/dT)^m\) to \((1-T)^{n+1}\) and evaluate at \(T=1\).
Every resulting term still vanishes there. Expanding gives the
alternating binomial identity, and linearity gives
\eqref{lm-finite-difference}. Its first coefficient is one, while
\(\sum_j\binom{n+1}{j}=2^{n+1}\) bounds each coefficient.
The rank statement follows from the literal identification
\(\mathcal O=\mathbb Z^2\), and every displayed congruence vanishes
on the kernel.
\end{proof}

For prime \(n>324\), the block \(B=n^4\), \(B\le k<2B\), contains
such consecutive indices. Thus a nonzero coefficient vector satisfying
arbitrarily many congruences need not have nonzero actual output.
This is a counterexample to that inference on the original block.
It does not exclude an argument that separately proves escape from
the kernel or retains a higher-dimensional actual output.

\begin{theorem}[An exact selector for different actual owners]
\label{lm-exact-selector}
Let \(Z_1,Z_2\in\mathcal O\) be primitive, with
\[
 \Delta=\det(Z_1,Z_2)\ne0,\qquad P(Z_1)P(Z_2)\ne0.
\]
For \(i=1,2\), let \(S_i\) be a nonempty finite set of primes
greater than three, with \(S_1\cap S_2=\varnothing\). For every
\(q\in S_i\) assume
\[
 e_q=v_q(P(Z_i))\ge1,\qquad q\nmid N(Z_i).
\]
Set
\[
 M_i=\prod_{q\in S_i}q^{e_q+1},\qquad M=M_1M_2.
\]
For every integer \(t\) satisfying
\begin{equation}
 t\equiv1\pmod {M_1},\qquad t\equiv0\pmod {M_2},
 \label{lm-selector-crt}
\end{equation}
define \(U_t=tZ_1+(1-t)Z_2\), \(g_t=\operatorname{cont}(U_t)\),
and \(V_t=U_t/g_t\). These elements are well-defined and
\begin{equation}
 g_t\mid |\Delta|,\qquad \gcd(g_t,M)=1,\qquad
 v_q(P(V_t))=e_q\quad(q\in S_1\cup S_2).
 \label{lm-content-and-depth}
\end{equation}
In particular the exact selected signed cost is
\begin{equation}
 \sum_{q\in S_1\cup S_2}\bigl(v_q(P(V_t))-3\bigr)\log q
   =\sum_{i=1}^2\sum_{q\in S_i}(e_q-3)\log q.
 \label{lm-selected-signed}
\end{equation}
There are infinitely many such integers \(t\).
\end{theorem}
\begin{proof}
The two moduli are coprime, so the integer Chinese remainder theorem
gives one residue class modulo \(M\). Linear independence of
\(Z_1,Z_2\) over \(\mathbb Q\) makes \(U_t\ne0\), since \(t\)
and \(1-t\) cannot both vanish. The actual determinant identities
\[
 \det(U_t,Z_2)=t\Delta,\qquad
 \det(U_t,Z_1)=-(1-t)\Delta
\]
imply \(g_t\mid\Delta\) by subtraction. For \(q\in S_i\),
\(U_t\equiv Z_i\pmod {q^{e_q+1}}\). The integer cubic therefore
has exactly valuation \(e_q\), and \(N(U_t)\) is a \(q\)-unit.
Consequently \(q\nmid g_t\), and dividing the cubic by \(g_t^3\)
preserves the exact valuation. This proves both asserted identities.
\end{proof}

There is no requirement that both input roots hit each marked prime.
The additional digit in \(q^{e_q+1}\) fixes the depth exactly,
including depths one and two. The theorem by itself makes no claim
about unmarked primes.

\begin{proposition}[The cost after actual primitive cancellation]
\label{lm-primitive-height}
For every selector in Theorem~\ref{lm-exact-selector},
\begin{align}
 |t|&\ge M_2,& |1-t|&\ge M_1,\notag\\
 \log|V_t|&\ge\log(\sqrt3/2)
       +\max\{\log M_2-\log|Z_2|,
                  \log M_1-\log|Z_1|\}.
 \label{lm-height-lower}
\end{align}
If \(n,B\) are positive integers and \(Z_i=w_{k_i}^n\) for integers
\(B\le k_i<2B\), then
\begin{equation}
 \log|V_t|\ge\tfrac12\log M-n\log(6B)+\log(\sqrt3/2).
 \label{lm-power-height-lower}
\end{equation}
The representative \(0\le t<M\) also satisfies
\begin{equation}
 |V_t|\le |U_t|<M(|Z_1|+|Z_2|).
 \label{lm-height-upper}
\end{equation}
\end{proposition}
\begin{proof}
The congruences exclude \(t=0,1\). Thus the nonzero integers \(t\)
and \(1-t\) are respectively multiples of \(M_2\) and \(M_1\).
The Gram determinant inequality in these coordinates is
\[
 |\det(X,Y)|\le \frac2{\sqrt3}|X|\,|Y|.
\]
Apply it to \((V_t,Z_2)\). Since \(|\Delta|/g_t\) is a positive
integer, it gives
\[
 \frac2{\sqrt3}|V_t|\,|Z_2|
     \ge |t|\frac{|\Delta|}{g_t}\ge |t|.
\]
The other determinant gives the analogous inequality for
\(|1-t|\) and \(Z_1\). Taking logarithms proves
\eqref{lm-height-lower}. The bounds \(|w_k|<6B\) and
\(\max(\log M_1,\log M_2)\ge\tfrac12\log M\) prove
\eqref{lm-power-height-lower}. For the chosen representative,
\(t\ne0,1\), \(|t|<M\), and \(|1-t|<M\); the triangle
inequality and \(g_t\ge1\) give \eqref{lm-height-upper}.
\end{proof}

The modulus cost is therefore present in an actual integer selector;
it does not disappear just by dividing out content. This is a bound
for this construction, not an optimality assertion for every map.

\begin{theorem}[A fixed-data positive-density family with full credit]
\label{lm-paid-affine-family}
Fix all the data of Theorem~\ref{lm-exact-selector}. Suppose also
\[
 Z_2=(a_2,b_2),\qquad Z_1-Z_2=(a_0,b_0),\qquad a_0,b_0>0.
\]
Choose a positive \(t_*\) in the CRT class with \(U_{t_*}\) having
positive coordinates. Let \(\mathcal E\) be the finite set of primes
dividing the nonzero integer
\[
 6M\Delta a_0b_0(a_0+b_0).
\]
For \(p\in\mathcal E\) put \(F_p=v_p(P(U_{t_*}))\), and define
\begin{gather}
 H=\prod_{p\in\mathcal E}p^{F_p+1},\qquad
 t(j)=t_*+Hj\quad(j\ge1),\qquad g=\operatorname{cont}(U_{t_*}),
 \notag\\
 f_p=F_p-3v_p(g),\qquad
 K=\prod_{p\in\mathcal E}p^{f_p}.
 \label{lm-frozen-data}
\end{gather}
Then \(M\mid H\), \(f_p\ge0\), and for every positive integer \(j\),
the actual primitive output
\[
 V_j=U_{t(j)}/g=(a_j,b_j)
\]
has positive coordinates. Its three arms are affine integer
polynomials in \(j\), with positive slopes.

There is a set \(\mathcal G\) of positive integers with natural density
\begin{equation}
 d(\mathcal G)=\prod_{p\notin\mathcal E}(1-3/p^2)\ge\tfrac14
 \label{lm-positive-density}
\end{equation}
such that, for every \(j\in\mathcal G\),
\begin{equation}
 T_j=a_jb_j(a_j+b_j)=K R_j,
 \quad R_j\text{ squarefree},\quad \gcd(K,R_j)=1,
 \label{lm-full-squarefree-support}
\end{equation}
and all primes of \(R_j\) lie outside \(\mathcal E\). The complete
signed identity and its limit are
\begin{align}
 J(V_j)&=-2\log T_j+3\log\frac{K}{\operatorname{rad}(K)},
 \label{lm-full-signed-identity}\\
 \frac{J(V_j)}{\log(a_j+b_j)}&\longrightarrow-6
       \quad(j\in\mathcal G,\ j\to\infty).
 \label{lm-signed-limit}
\end{align}
For a constant \(C\) depending on these fixed data,
\(a_j+b_j\le C\operatorname{rad}(T_j)^{1/3}\) for every
\(j\in\mathcal G\).
\end{theorem}
\begin{proof}
At a marked prime, the selector theorem gives \(F_q=e_q\), so
\(M\mid H\). For any \(p\in\mathcal E\), the congruence
\(U_{t(j)}\equiv U_{t_*}\pmod {p^{F_p+1}}\) preserves the exact
cubic valuation. It also preserves the common coordinate valuation:
if \(v=\min(v_p(a_*),v_p(b_*))\), then \(F_p\ge3v\), whence
\(F_p+1>v\), and a coordinate of valuation \(v\) retains it.
No prime outside \(\mathcal E\) divides the content, by
\eqref{lm-content-and-depth}. Thus \(\operatorname{cont}(U_{t(j)})=g\)
for every \(j\). In particular \(g\mid H\), so the quotient arms
are affine integer polynomials, with positive slopes and positive
initial coordinates. Dividing the cubic by \(g^3\) gives the
nonnegative exact depths \(f_p\).

For \(p\notin\mathcal E\), all three slopes are \(p\)-units.
The determinant of the slope and constant coordinate pairs is
\(H\Delta/g^2\), up to sign, and is also a \(p\)-unit. Each
of the three affine arms therefore has one simple root modulo
\(p^2\), and their roots are distinct already modulo \(p\).
Exactly three classes modulo \(p^2\) are forbidden by
\(p^2\mid T_j\). Notice that every such prime is at least five.

For any fixed cutoff \(Y\), finite CRT gives density
\(\prod_{p\le Y,\ p\notin\mathcal E}(1-3/p^2)\) for avoiding
these finitely many classes. The entire omitted tail is elementary.
There is a fixed \(C_0>0\) such that every individual positive
affine arm at \(1\le j\le X\) is at most \(C_0X\), for every
real \(X\ge1\). Distinct roots modulo \(p\) ensure that
\(p^2\mid T_j\) forces \(p^2\) to divide one arm, rather than
arising from two different arms. Such a prime is at most
\(\sqrt{C_0X}\). Counting the three residue classes gives
\begin{equation}
 \#\{1\le j\le X:\exists p>Y,\ p\notin\mathcal E,
                         \ p^2\mid T_j\}
 \le3X\sum_{p>Y}p^{-2}+3\sqrt{C_0X}.
 \label{lm-complete-square-tail}
\end{equation}
Divide by \(X\), let \(X\to\infty\), and then let \(Y\to\infty\).
The finite CRT upper and lower bounds prove existence of the natural
density and its exact infinite product. Its lower bound follows from
\[
 \sum_{p\ge5}\frac3{p^2}
 \le3\sum_{m\ge5}\frac1{m^2}
 <3\int_4^\infty x^{-2}\,dx=\tfrac34
\]
and the finite inequality \(\prod(1-u_i)\ge1-\sum u_i\) for
\(0\le u_i\le1\). Passing to the convergent product proves
\eqref{lm-positive-density}; removing additional primes in
\(\mathcal E\) only increases it. No distribution theorem for
primes is needed.

At each prime of \(\mathcal E\) the exact depth is \(f_p\).
Outside it a good parameter has depth zero or one. This proves
\eqref{lm-full-squarefree-support}, and hence
\[
 \operatorname{rad}(T_j)=\operatorname{rad}(K)R_j
      =\frac{T_j}{K/\operatorname{rad}(K)}.
\]
Substitution in \eqref{lm-signed-cost} gives
\eqref{lm-full-signed-identity}. The positive affine slopes imply
\[
 \log T_j=3\log j+O(1),\qquad
 \log(a_j+b_j)=\log j+O(1),
\]
which prove \eqref{lm-signed-limit}. They also give
\(T_j\ge c_0(a_j+b_j)^3\) for a fixed \(c_0>0\), decreasing
\(c_0\) if necessary for finitely many small indices. The radical
identity proves the last bound.
\end{proof}

The data, progression and threshold in this theorem are fixed. The
numbers \(K\) and \(C_0\) can be very large when the input roots or
packets change. The density lower bound is a natural-density statement
in this fixed progression; it is not a density statement in a moving
original root block.

\begin{proposition}[Actual source roots and nonempty examples]
\label{lm-actual-source-roots}
Let \(n\ge2\), \(B\ge n\) be integers, and let
\(B\le k_2<k_1<2B\). The two actual powers
\(Z_i=(3k_i+\zeta)^n\) are primitive and unramified, have
\(P(Z_i)>0\), are nonparallel, and have both coordinates of
\(Z_1-Z_2\) positive.
\end{proposition}
\begin{proof}
The element \(3k+\zeta\) is primitive and its norm is one modulo
three. Its powers remain primitive: modulo any prime other than
three the Eisenstein algebra is a product of fields or a field, so
an element with zero power is zero; the coordinate one prevents that.
At three the norm is a unit. This proves the required property at
every prime and also nonramification.

For real \(x\ge B\), put
\[
 \theta(x)=\arg(3x+\zeta)
       =\arctan\frac{\sqrt3}{6x+1}.
\]
This function is positive, strictly decreasing, and less than
\(1/(3x)\). Thus \(0<n\theta(x)<1/3<\pi/3\), so each power
lies in the open positive coordinate cone. Their two arguments
are distinct and differ by less than \(\pi\), proving
nonparallelism. Finally
\[
 (3k_1+\zeta)^n-(3k_2+\zeta)^n
    =\int_{k_2}^{k_1}3n(3x+\zeta)^{n-1}\,dx.
\]
Since \(n\ge2\), the integrand lies strictly in the same cone;
both coordinates of its integral are positive.
\end{proof}

For \(n=2\), \(B=2\), the roots \(w_2=6+\zeta\),
\(w_3=9+\zeta\) give
\[
 Z_2=(35,13),\quad P(Z_2)=35\cdot13\cdot48,
 \qquad Z_1=(80,19),\quad P(Z_1)=80\cdot19\cdot99.
\]
The prime 19 belongs only to the first boundary, and 13 only to the
second. Their depths are one and their selected norms are units.
The choices \(S_1=\{19\}\), \(S_2=\{13\}\) therefore give a
nonempty infinite actual output family in
Theorem~\ref{lm-paid-affine-family}.

Positive signed costs can also be exhibited with different owners.
Still take \(n=2\), and set
\[
 k_1=21720+2\cdot19^5=4973918,\qquad
 k_2=61882+13^6=4888691.
\]
Both belong to the extended block \(B=4880000\). Their second
coordinates are respectively \(19^4\cdot229\) and
\(13^5\cdot79\). If \(q>3\) divides the second coordinate
\(6k+1\), the first coordinate \((3k)^2-1\) is
\(-3/4\pmod q\). Thus the actual whole boundary depths are
four and five. The other owner's boundary is a unit at the marked
prime, as follows from \(k_1\equiv1\pmod {13}\) and
\(k_2\equiv10\pmod {19}\). The selected positive signed cost
is exactly \(\log19+2\log13\). The selector preserves it, while
the new squarefree support in the proved affine family pays it.

These examples have exponent two. They belong to the extended
actual-power family of Proposition~\ref{lm-actual-source-roots};
they do not satisfy the original restrictions of prime \(n>324\),
\(B=n^4\), or large marked primes beyond that block. The selector
lemma applies to roots in that original block when the required
actual packets are supplied, but these examples do not establish
such moving-block membership.

\paragraph{Evidence and the remaining arithmetic constraint.}
The standard-library script
\texttt{replay\_affine\_\allowbreak selectors.py --check} checks ten exact CRT
selectors, including the depth-four and depth-five packet, 256 actual
fixed-progression content/depth rows, and the three distinct simple
roots modulo \(p^2\) for 162 eligible primes up to 997. It also checks
six literal finite-difference identities at exponents
\(2,5,7,11,17,331\); the last is in the original block \(B=n^4\).
The canonical certificate SHA256 is
\begin{quote}\ttfamily
70945a8da8fdecee\allowbreak 2a8857dd9ea71ab2\allowbreak
bb12f9054a8f25b3\allowbreak 5ad87f342df6d767.
\end{quote}
These are finite integer checks. They do not factor every unmarked
arm or prove density by sampling. The complete infinite tail and
density proof is \eqref{lm-complete-square-tail}, with the finite
prime exceptions retained in \(K\). There is no new Lean claim here.

The positive theorem is about new primitive affine outputs and their
full signed costs. It neither bounds the original inputs' tails nor
proves that the outputs are powers or have small-residual power
representations. Its constants cannot be applied uniformly to moving
packets. The additional constraint would be to retain useful original
compression while proving a sufficient supply of new negative credit.
Original boundary tails, exceptional roots, the complement of the
independent domain, and general ABC coverage remain open.
